% !TEX root =esercizi.tex
\chapter{Lezioni 4--6 (30 settembre -- 5 ottobre 2026)}

\begin{vuoto}[Argomenti svolti]{}
\begin{itemize}
\item[\bul] \textbf{Lezione 4 --- Equazioni e disequazioni di primo grado}. Principi di equivalenza; equazioni di primo grado e fratte, condizioni di esistenza. Disequazioni di primo grado; sistemi di disequazioni. Equazioni e disequazioni con il valore assoluto.
\item[\bul] \textbf{Lezione 5 --- Equazioni e disequazioni di secondo grado}. Prodotti notevoli e scomposizioni; legge di annullamento del prodotto. Formula risolutiva, discriminante, scomposizione del trinomio. Segno del trinomio e metodo della parabola. Disequazioni di secondo grado; segno di prodotti e quozienti (metodo dei due sistemi, tabella dei segni), anche con fattori di primo grado (Cap.~4, \S4.5); biquadratiche, sistemi. Notazione scientifica (Cap.~6, \S6.7) e percentuali (Cap.~7, \S7.1).
\item[\bul] \textbf{Lezione 6 --- Completezza, potenze e radici}. Maggioranti e minoranti, estremo superiore e inferiore, massimo e minimo; proprietà di Archimede; completezza di $\R$. Potenze a esponente intero e razionale, cenno all'esponente reale ($3^{\sqrt2}$ come estremo superiore); radici $n$-esime, $\sqrt{x^2}=|x|$. Disequazioni $x^n\gtrless a$.
\end{itemize}
\end{vuoto}

Unità di misura, ordini di grandezza ed equazioni e disequazioni con la radice quadrata non sono ancora stati svolti: i relativi esercizi sono rinviati alla Scheda~3. La caratterizzazione dell'estremo superiore e inferiore con $\eps$, che sarà svolta in L7, è richiamata nel riquadro all'inizio della sezione~2.8.

Il livello di difficoltà è indicato con \liv{1} (esercizio di base), \liv{2} (esercizio standard), \liv{3} (esercizio più impegnativo). In tutti gli esercizi la soluzione va scritta come intervallo o unione di intervalli.

%-------------------------------------------------------------------
\section{Equazioni e disequazioni di primo grado}

\begin{esercizio}[\liv{1}]
Risolvere: \ $2x+3=7$; \ $5-x=2x-4$; \ $\dfrac x3+1=\dfrac x2$.
\end{esercizio}
\risp{$x=2$; \ $x=3$; \ $1=\frac x6$, cioè $x=6$.}

\begin{esercizio}[\liv{1}]
Tradurre in un'equazione e risolvere: ``il doppio di un numero, diminuito di $5$, è uguale a $11$''.
\end{esercizio}
\risp{$2x-5=11$, $x=8$.}

\begin{esercizio}[\liv{1}]
Risolvere, dicendo se l'equazione è determinata, impossibile o indeterminata:
\[
3(x-2)-2(x+1)=x-8,\qquad \frac{x+1}{2}-\frac{x-1}{3}=1,\qquad 2(x-1)=2x+3 .
\]
\end{esercizio}
\risp{La prima si riduce a $x-8=x-8$: indeterminata. La seconda dà $3x+3-2x+2=6$, cioè $x=1$. La terza dà $-2=3$: impossibile.}

\begin{esercizio}[\liv{1}]
Risolvere e rappresentare sulla retta:
\[
5-2x\ge3x+20,\qquad \frac{x-2}{4}>\frac{2x+1}{3},\qquad 0{,}2x-1{,}5<0{,}5x .
\]
\end{esercizio}
\risp{$-5x\ge15$, cioè $x\le-3$; \ $3x-6>8x+4$, cioè $-5x>10$, cioè $x<-2$; \ $-1{,}5<0{,}3x$, cioè $x>-5$.}

\begin{esercizio}[\liv{2}]
Risolvere, scrivendo prima le condizioni di esistenza:
\[
\frac{3}{x-1}=\frac{2}{x+2},\qquad \frac{x}{x-2}-\frac{2}{x-2}=3 .
\]
\end{esercizio}
\risp{Prima: $x\ne1,-2$; $3x+6=2x-2$, $x=-8$, accettabile. Seconda: $x\ne2$; per $x\ne2$ il primo membro vale $\frac{x-2}{x-2}=1\ne3$: impossibile.}

\begin{esercizio}[\liv{2}]
Un bacino di raccolta contiene $12\,000$ L; da esso si prelevano $150$ L all'ora e vi affluiscono $90$ L all'ora. Dopo quante ore il contenuto scende sotto i $3000$ L? In quale intervallo di tempo il contenuto è compreso fra $6000$ e $9000$ L?
\end{esercizio}
\risp{Il contenuto è $12\,000-60t$ (con $t$ in ore). $12\,000-60t<3000$ dà $t>150$ h. $6000\le12\,000-60t\le9000$ dà $50\le t\le100$ h.}

%-------------------------------------------------------------------
\section{Sistemi di disequazioni}

\begin{esercizio}[\liv{1}]
Risolvere e rappresentare sulla retta:
\[
\begin{cases}x-1>0\\ x<4\end{cases}\qquad \begin{cases}2x\ge6\\ x-5\le0\end{cases}\qquad \begin{cases}x>2\\ x<-1\end{cases}
\]
\end{esercizio}
\risp{$]1,4[$; \ $[3,5]$; \ nessuna soluzione (le due semirette non si intersecano).}

\begin{esercizio}[\liv{1}]
Risolvere:
\[
\begin{cases}2x+1>3\\ 4-x\ge0\end{cases}\qquad
\begin{cases}3x-2\le x\\ x+5>2x\end{cases}\qquad
-1\le\frac{3-x}{2}<2 .
\]
\end{esercizio}
\risp{$]1,4]$; \ $x\le1$ e $x<5$, cioè $]-\infty,1]$; \ $-2\le3-x<4$, cioè $-1<x\le5$: $]-1,5]$.}

\begin{esercizio}[\liv{2}]
Risolvere il sistema $\begin{cases}\frac{x-3}{2}\ge x-5\\ 2(x+1)>3x-4\\ x\ge0\end{cases}$.
\end{esercizio}
\risp{La prima dà $x\le7$, la seconda $x<6$: con $x\ge0$ la soluzione è $[0,6[$.}

%-------------------------------------------------------------------
\section{Segno di prodotti e quozienti}

\begin{esercizio}[\liv{1}]
Per quali $x$ il quoziente $\dfrac{x-3}{x+2}$ è positivo? Negativo? Nullo? Per quali non è definito?
\end{esercizio}
\risp{Positivo per $x<-2$ oppure $x>3$; negativo per $-2<x<3$; nullo per $x=3$; non definito per $x=-2$.}

\begin{esercizio}[\liv{1}]
Risolvere $(x-1)(x-4)>0$ e $(x-1)(x-4)\le0$.
\end{esercizio}
\risp{$]-\infty,1[\,\cup\,]4,+\infty[$; \ $[1,4]$.}

\begin{esercizio}[\liv{1}]
Risolvere:
\[
\frac{x+3}{2-x}\ge0,\qquad (2x-1)(x+4)<0,\qquad \frac{3-x}{x+1}\le0 .
\]
\end{esercizio}
\risp{$[-3,2[$; \ $]-4,\frac12[$; \ $]-\infty,-1[\,\cup\,[3,+\infty[$.}

\begin{esercizio}[\liv{2}]
Risolvere, con la tabella dei segni:
\[
\frac{(x-1)(x+2)}{x-4}\ge0,\qquad \frac{2}{x-1}\ge1,\qquad \frac1{x-1}<\frac1{x+1} .
\]
\end{esercizio}
\risp{Prima: $[-2,1]\,\cup\,]4,+\infty[$. Seconda: $\frac{3-x}{x-1}\ge0$, cioè $]1,3]$. Terza: $\frac{2}{(x-1)(x+1)}<0$, cioè $(x-1)(x+1)<0$: $]-1,1[$.}

\begin{esercizio}[\liv{3}]
Si mescolano $x$ litri di una soluzione al $30\%$ con $5$ litri di una soluzione al $10\%$. Scrivere la concentrazione della miscela in funzione di $x$ e determinare per quali $x$ essa supera il $20\%$. Perché in questo caso si può moltiplicare per il denominatore senza studiarne il segno?
\end{esercizio}
\risp{$c(x)=\frac{0{,}3x+0{,}5}{x+5}$. Per $x\ge0$ il denominatore è positivo, quindi $c>0{,}2$ equivale a $0{,}3x+0{,}5>0{,}2x+1$, cioè $x>5$ L.}

%-------------------------------------------------------------------
\section{Valore assoluto}

\begin{esercizio}[\liv{1}]
Calcolare $|x-3|$ per $x=5$, $x=3$, $x=-1$. Scrivere poi $|x-3|$ senza il valore assoluto, distinguendo i casi $x\ge3$ e $x<3$.
\end{esercizio}
\risp{$2$, $0$, $4$. \ $|x-3|=x-3$ se $x\ge3$, \ $|x-3|=3-x$ se $x<3$.}

\begin{esercizio}[\liv{1}]
Risolvere: \ $|x|=4$; \ $|x|<2$; \ $|x|\ge1$; \ $|x|=-3$.
\end{esercizio}
\risp{$x=\pm4$; \ $]-2,2[$; \ $]-\infty,-1]\cup[1,+\infty[$; \ nessuna soluzione.}

\begin{esercizio}[\liv{1}]
Risolvere:
\begin{gather*}
|2x-3|=5,\qquad |x+1|<3,\\
|4-x|\ge1,\qquad |3x+2|\le-1 .
\end{gather*}
\end{esercizio}
\risp{$x=4$ oppure $x=-1$; \ $]-4,2[$; \ $]-\infty,3]\cup[5,+\infty[$; \ nessuna soluzione.}

\begin{esercizio}[\liv{2}]
Risolvere:
\begin{gather*}
|x-2|=|x+4|,\qquad |x|+|x-2|=4,\\
|x-1|>2x,\qquad |x-1|+|x+1|<4 .
\end{gather*}
\end{esercizio}
\risp{Prima: punto equidistante da $2$ e $-4$, $x=-1$. Seconda: per $x<0$, $2-2x=4$, $x=-1$; per $0\le x<2$, $2=4$, mai; per $x\ge2$, $2x-2=4$, $x=3$; soluzioni $-1$ e $3$. Terza: per $x\ge1$, $x-1>2x$ dà $x<-1$, incompatibile; per $x<1$, $1-x>2x$ dà $x<\frac13$; soluzione $]-\infty,\frac13[$. Quarta: $]-2,2[$.}

\begin{esercizio}[\liv{3}]
Una specie di pesci vive fra $8$ e $24\,^\circ$C. In un lago, la temperatura alla profondità $h$ (in metri, $0\le h\le20$) è $T(h)=22-0{,}8h$. Scrivere la condizione sulla temperatura con un valore assoluto e determinare a quali profondità la specie trova condizioni adatte.
\end{esercizio}
\risp{$|T-16|\le8$, cioè $|6-0{,}8h|\le8$: $-8\le6-0{,}8h\le8$, cioè $-2{,}5\le h\le17{,}5$. Con $0\le h\le20$: fra $0$ e $17{,}5$ m.}

%-------------------------------------------------------------------
\section{Equazioni di secondo grado e scomposizioni}

\begin{esercizio}[\liv{1}]
Scomporre in fattori:
\begin{gather*}
x^2-10x+25,\qquad 4x^2-1,\qquad x^3+27,\\
x^3-x^2-4x+4,\qquad 2x^2-7x+3 .
\end{gather*}
\end{esercizio}
\risp{$(x-5)^2$; \ $(2x-1)(2x+1)$; \ $(x+3)(x^2-3x+9)$; \ $(x-1)(x-2)(x+2)$; \ $(x-3)(2x-1)$.}

\begin{esercizio}[\liv{1}]
Risolvere, usando la formula risolutiva solo quando serve:
\begin{gather*}
x^2-x-12=0,\qquad 5x^2-20x=0,\qquad 9x^2-4=0,\\
x^2+6x+10=0,\qquad x^2-2\sqrt2\,x+2=0 .
\end{gather*}
\end{esercizio}
\risp{$4$ e $-3$; \ $0$ e $4$; \ $\pm\frac23$; \ nessuna soluzione ($\Delta=-4$); \ $(x-\sqrt2)^2=0$, cioè $x=\sqrt2$.}

\begin{esercizio}[\liv{2}]
(a) Risolvere $(x-1)(x+2)=4$ e $\dfrac{x}{x-1}+\dfrac1x=2$. (b) Per quali $k$ l'equazione $x^2-4x+k=0$ ha due soluzioni distinte? (c) Trovare due numeri di somma $10$ e prodotto $21$.
\end{esercizio}
\risp{(a) $x^2+x-6=0$: $2$ e $-3$. Con $x\ne0,1$: $x^2+x-1=2x^2-2x$, cioè $x^2-3x+1=0$, $x=\frac{3\pm\sqrt5}2$, entrambe accettabili. (b) $\Delta=16-4k>0$: $k<4$. (c) Sono le soluzioni di $t^2-10t+21=0$: $3$ e $7$.}

%-------------------------------------------------------------------
\section{Disequazioni di secondo grado}

\begin{esercizio}[\liv{1}]
Risolvere: \ $x^2<9$; \ $x^2\ge4$; \ $x^2-x\ge0$; \ $x^2+4>0$.
\end{esercizio}
\risp{$]-3,3[$; \ $]-\infty,-2]\cup[2,+\infty[$; \ $x(x-1)\ge0$: $]-\infty,0]\cup[1,+\infty[$; \ ogni $x\in\R$.}

\begin{esercizio}[\liv{1}]
Risolvere:
\begin{gather*}
x^2-3x-10\le0,\qquad 2x^2-x-1>0,\qquad -x^2+2x-1\ge0,\\
x^2+x+1>0,\qquad 4-x^2<0 .
\end{gather*}
\end{esercizio}
\risp{$[-2,5]$; \ $]-\infty,-\frac12[\,\cup\,]1,+\infty[$; \ $-(x-1)^2\ge0$, cioè $x=1$; \ ogni $x\in\R$; \ $]-\infty,-2[\,\cup\,]2,+\infty[$.}

\begin{esercizio}[\liv{2}]
Risolvere:
\begin{gather*}
\frac{x^2-4}{x+1}\ge0,\qquad \frac{x^2-5x+6}{x^2+x+1}<0,\\
(x-2)(x^2-2x-3)>0,\qquad \frac{x}{x-2}\ge\frac{3}{x+2} .
\end{gather*}
\end{esercizio}
\risp{Prima: $[-2,-1[\,\cup\,[2,+\infty[$. Seconda: il denominatore è sempre positivo: $]2,3[$. Terza: $(x-2)(x-3)(x+1)>0$: $]-1,2[\,\cup\,]3,+\infty[$. Quarta: $\frac{x^2-x+6}{(x-2)(x+2)}\ge0$, con numeratore sempre positivo: $]-\infty,-2[\,\cup\,]2,+\infty[$.}

\begin{esercizio}[\liv{2}]
Risolvere il sistema $\begin{cases}x^2-5x+4\le0\\ x^2-4>0\end{cases}$ e la disequazione biquadratica $x^4-13x^2+36\le0$.
\end{esercizio}
\risp{Sistema: $[1,4]$ intersecato con $]-\infty,-2[\,\cup\,]2,+\infty[$: $]2,4]$. Biquadratica: con $t=x^2$, $4\le t\le9$, cioè $2\le|x|\le3$: $[-3,-2]\cup[2,3]$.}

\begin{esercizio}[\liv{3}]
Un prato rettangolare ha perimetro $200$ m. Per quali lunghezze di un lato l'area supera $2100\ \mathrm{m^2}$? Qual è l'area massima possibile? (Completare il quadrato.)
\end{esercizio}
\risp{Con lati $x$ e $100-x$: $x(100-x)>2100$, cioè $x^2-100x+2100<0$, cioè $30<x<70$ m. Poiché $x(100-x)=2500-(x-50)^2$, l'area massima è $2500\ \mathrm{m^2}$, per il quadrato di lato $50$ m.}

%-------------------------------------------------------------------
\section{Valore assoluto e radici}

\begin{esercizio}[\liv{1}]
Scrivere le condizioni di esistenza di $\sqrt{3-x}$, \ $\sqrt{x^2-1}$, \ $\dfrac{1}{\sqrt{x+4}}$.
\end{esercizio}
\risp{$x\le3$; \ $x^2\ge1$, cioè $x\le-1$ oppure $x\ge1$; \ $x>-4$ (la radice deve esistere ed essere diversa da zero).}

\begin{esercizio}[\liv{1}]
Risolvere: \ $\sqrt{x-2}=3$; \ $\sqrt{x+1}=-2$; \ $\sqrt{x^2}=5$.
\end{esercizio}
\risp{$x-2=9$, $x=11$; \ nessuna soluzione (una radice non è mai negativa); \ $|x|=5$, cioè $x=\pm5$.}

\begin{esercizio}[\liv{2}]
Risolvere:
\[
|x^2-4|<3,\qquad |x^2-x|=2,\qquad x^2-5|x|+6<0 .
\]
\end{esercizio}
\risp{Prima: $1<x^2<7$, cioè $1<|x|<\sqrt7$: $]-\sqrt7,-1[\,\cup\,]1,\sqrt7[$. Seconda: $x^2-x=2$ dà $2$ e $-1$; $x^2-x=-2$ non ha soluzioni. Terza: con $t=|x|$, $2<t<3$: $]-3,-2[\,\cup\,]2,3[$.}

%-------------------------------------------------------------------
\section{Estremo superiore e completezza}

\begin{vuoto}[Caratterizzazione dell'estremo superiore e inferiore]{}
Sia $A\subseteq\R$ non vuoto. Un numero $s$ è l'estremo superiore di $A$ se e solo se
\begin{itemize}
\item[(i)] $x\le s$ per ogni $x\in A$ \hfill ($s$ è un maggiorante);
\item[(ii)] per ogni $\eps>0$ esiste $x\in A$ con $x>s-\eps$ \hfill ($s-\eps$ non è un maggiorante).
\end{itemize}
Analogamente, $i$ è l'estremo inferiore di $A$ se e solo se
\begin{itemize}
\item[(i)] $x\ge i$ per ogni $x\in A$ \hfill ($i$ è un minorante);
\item[(ii)] per ogni $\eps>0$ esiste $x\in A$ con $x<i+\eps$ \hfill ($i+\eps$ non è un minorante).
\end{itemize}
Nella (ii) l'ordine dei quantificatori conta: prima si fissa $\eps$, poi si cerca $x$, che in generale dipende da $\eps$. Negli esempi con $\frac1n$ la (ii) si verifica con la \emph{proprietà di Archimede}: per ogni $\eps>0$ esiste $n\in\N$ con $n>\frac1\eps$, cioè $\frac1n<\eps$.
\end{vuoto}

\begin{esercizio}[\liv{1}]
Per ciascun insieme dire se ha massimo, minimo, estremo superiore, estremo inferiore: \ $[2,5[$; \ $]0,+\infty[$; \ $\{1,2,3\}$.
\end{esercizio}
\risp{$[2,5[$: minimo $2$, estremo superiore $5$, nessun massimo. \ $]0,+\infty[$: estremo inferiore $0$, nessun minimo; non è limitato superiormente, quindi non ha estremo superiore in $\R$. \ $\{1,2,3\}$: minimo $1$, massimo $3$.}

\begin{esercizio}[\liv{1}]
Per ciascun insieme determinare maggioranti, estremo superiore e inferiore, e dire se esistono massimo e minimo:
\begin{gather*}
]-1,3],\qquad \Bigl\{\frac1n\colon n\in\N\Bigr\},\\
\{x\in\R\colon x^2\le4\},\qquad \Bigl\{(-1)^n+\frac1n\colon n\in\N\Bigr\}.
\end{gather*}
\end{esercizio}
\risp{$]-1,3]$: $\sup=\max=3$, $\inf=-1$, non minimo. $\{1/n\}$: $\sup=\max=1$, $\inf=0$, non minimo. $[-2,2]$: $\max=2$, $\min=-2$. Nell'ultimo i termini sono $0,\frac32,-\frac23,\frac54,\ldots$: $\sup=\max=\frac32$ (per $n=2$), $\inf=-1$, non minimo.}

\begin{esercizio}[\liv{1}]
Determinare, se esistono, massimo, minimo, estremo superiore ed estremo inferiore di:
\begin{gather*}
\N,\qquad \Z,\qquad \Bigl\{-\frac1n\colon n\in\N\Bigr\},\\
\Bigl\{\frac{n}{n+1}\colon n\in\N\Bigr\},\qquad \{x\in\R\colon |x-2|<1\},\qquad ]-\infty,3]\cup\{5\}.
\end{gather*}
\end{esercizio}
\risp{$\N$: $\min=1$, non limitato superiormente ($\sup\N=+\infty$, per Archimede). \ $\Z$: $\sup=+\infty$, $\inf=-\infty$. \ $\{-1,-\frac12,-\frac13,\ldots\}$: $\min=-1$, $\sup=0$, non massimo. \ $\{\frac12,\frac23,\frac34,\ldots\}$: $\min=\frac12$, $\sup=1$, non massimo. \ $]1,3[$: $\inf=1$, $\sup=3$, né minimo né massimo. \ $\max=5$, $\inf=-\infty$.}

\begin{esercizio}[\liv{1}]
Vero o falso?
(a) Ogni insieme finito non vuoto ha massimo.
(b) Se $A$ ha massimo, allora $\max A=\sup A$.
(c) Se $\sup A\in A$, allora $\sup A$ è il massimo di $A$.
(d) Ogni insieme non vuoto e limitato superiormente ha massimo.
(e) Se $M$ è un maggiorante di $A$, anche $M+1$ lo è.
(f) Un insieme può avere due estremi superiori diversi.
\end{esercizio}
\risp{(a) V. (b) V: il massimo è un maggiorante, e ogni altro maggiorante è $\ge$ di esso perché il massimo sta in $A$. (c) V. (d) F: $]0,1[$. (e) V. (f) F: il minimo dei maggioranti, se esiste, è uno solo.}

\begin{esercizio}[\liv{1}]
Una centralina registra la temperatura dell'aria ogni ora, per un mese. (a) L'insieme delle temperature registrate ha sempre massimo e minimo? Perché? (b) L'insieme $\{15-\frac1n\colon n\in\N\}$ descrive temperature che si avvicinano a $15\,^\circ$C senza raggiungerlo: determinarne massimo, minimo, estremo superiore ed estremo inferiore.
\end{esercizio}
\risp{(a) Sì: è un insieme finito (al più $31\cdot24=744$ valori), e fra un numero finito di numeri ce n'è sempre uno più grande e uno più piccolo di tutti. (b) $\min=14$ (per $n=1$), $\sup=15$, non massimo.}

\begin{esercizio}[\liv{2}]
Usando la caratterizzazione dell'estremo superiore, provare che $\sup\{1-\frac1n\colon n\in\N\}=1$. Quale proprietà di $\R$ si usa?
\end{esercizio}
\risp{(i) $1-\frac1n\le1$ per ogni $n$. (ii) Dato $\eps>0$, per la proprietà di Archimede esiste $n>\frac1\eps$, e allora $1-\frac1n>1-\eps$.}

\begin{esercizio}[\liv{2}]
Usando la caratterizzazione dell'estremo inferiore, provare che $\inf\{2+\frac3n\colon n\in\N\}=2$. L'insieme ha minimo? Ha massimo?
\end{esercizio}
\risp{(i) $2+\frac3n>2$ per ogni $n$. (ii) Dato $\eps>0$, per Archimede esiste $n>\frac3\eps$, e allora $2+\frac3n<2+\eps$. Non ha minimo ($2$ non è della forma $2+\frac3n$); ha massimo $5$, per $n=1$.}

\begin{esercizio}[\liv{2}]
Determinare estremo superiore ed estremo inferiore di $\{x^2\colon -1<x<2\}$ e di $\{\frac1x\colon x>0\}$, dicendo se sono massimo e minimo. (Disegnare il grafico di $x^2$ aiuta.)
\end{esercizio}
\risp{Al variare di $x$ in $]-1,2[$, $x^2$ descrive $[0,4[$: $\min=0$ (per $x=0$), $\sup=4$, non massimo. Al variare di $x>0$, $\frac1x$ descrive $]0,+\infty[$: $\inf=0$, non minimo; non limitato superiormente.}

\begin{esercizio}[\liv{2}]
Siano $A\subseteq B$ due insiemi non vuoti, con $B$ limitato superiormente. Provare che $A$ è limitato superiormente e che $\sup A\le\sup B$. Mostrare con un esempio che può essere $\sup A=\sup B$ anche se $A\ne B$.
\end{esercizio}
\risp{$\sup B$ è un maggiorante di $B$, quindi anche di $A$ (ogni elemento di $A$ sta in $B$). Allora $A$ è limitato superiormente, e $\sup A$, che è il più piccolo dei maggioranti di $A$, è $\le\sup B$. Esempio: $A=]0,1[$, $B=[0,1]$.}

\begin{esercizio}[\liv{2}]
Sia $A=\{x\in\Q\colon x>0,\ x^2<3\}$. Qual è l'estremo superiore di $A$ in $\R$? Appartiene a $\Q$? Che cosa dice questo esempio su $\Q$?
\end{esercizio}
\risp{$\sup A=\sqrt3$, che è irrazionale (esercizio della Scheda~1): come sottoinsieme di $\Q$, $A$ è limitato ma non ha estremo superiore in $\Q$. È l'esempio che mostra che $\Q$ non è completo.}

\begin{esercizio}[\liv{3}]
Sia $A$ non vuoto e limitato superiormente, e sia $-A=\{-x\colon x\in A\}$. Provare che $-A$ è limitato inferiormente e che $\inf(-A)=-\sup A$.
\end{esercizio}
\risp{Sia $s=\sup A$. (i) $x\le s$ per ogni $x\in A$ equivale a $-x\ge-s$: $-s$ è un minorante di $-A$. (ii) Dato $\eps>0$ esiste $x\in A$ con $x>s-\eps$, cioè $-x<-s+\eps$, con $-x\in-A$. Per la caratterizzazione dell'estremo inferiore, $\inf(-A)=-s$.}

%-------------------------------------------------------------------
\section{Potenze e radici}

\begin{esercizio}[\liv{1}]
Calcolare senza calcolatrice: \ $2^{-3}$, \ $\bigl(\tfrac12\bigr)^{-2}$, \ $10^{0}$, \ $4^{1/2}$, \ $8^{2/3}$.
\end{esercizio}
\risp{$\frac18$; \ $4$; \ $1$; \ $2$; \ $4$.}

\begin{esercizio}[\liv{1}]
Semplificare: \ $a^3a^5$, \ $(a^2)^4$, \ $\dfrac{a^7}{a^3}$ (con $a\ne0$), \ $(2a)^3$.
\end{esercizio}
\risp{$a^8$; \ $a^8$; \ $a^4$; \ $8a^3$.}

\begin{esercizio}[\liv{1}]
Calcolare:
\begin{gather*}
\bigl(2^3\cdot2^{-5}\bigr)^{-1},\qquad 27^{2/3},\qquad 16^{-3/4},\\
\sqrt{50}-\sqrt{18},\qquad \sqrt[3]{-64} .
\end{gather*}
\end{esercizio}
\risp{$4$; \ $9$; \ $\frac18$; \ $5\sqrt2-3\sqrt2=2\sqrt2$; \ $-4$.}

\begin{esercizio}[\liv{2}]
Semplificare, indicando le condizioni:
\begin{gather*}
\sqrt{x^2-2x+1},\qquad \sqrt[3]{(x-2)^3},\\
\sqrt{x^4},\qquad \bigl(x^{1/2}x^{1/3}\bigr)^6 .
\end{gather*}
\end{esercizio}
\risp{$|x-1|$; \ $x-2$; \ $x^2$; \ $x^5$, con $x>0$.}

\begin{esercizio}[\liv{2}]
Risolvere:
\begin{gather*}
x^3=125,\quad x^4=81,\quad x^6=-1,\quad x^3\le-8,\\
x^4\ge16,\quad x^5+1>0,\quad \sqrt[3]{2x-1}<1,\quad \sqrt[4]{x+2}\ge2 .
\end{gather*}
\end{esercizio}
\risp{$5$; \ $\pm3$; \ nessuna; \ $x\le-2$; \ $|x|\ge2$; \ $x>-1$; \ $2x-1<1$, $x<1$; \ $x+2\ge16$, $x\ge14$.}

\begin{esercizio}[\liv{3}]
Confrontare senza calcolatrice $\sqrt[3]{3}$ e $\sqrt2$; poi ordinare $2^{1/2}$, $3^{1/3}$, $5^{1/5}$.
\end{esercizio}
\risp{Elevando alla sesta (potenza crescente sui positivi): $9>8$, quindi $\sqrt[3]3>\sqrt2$. Elevando alla trentesima: $2^{15}=32\,768$, $3^{10}=59\,049$, $5^6=15\,625$, quindi $5^{1/5}<2^{1/2}<3^{1/3}$.}

%-------------------------------------------------------------------
\section{Notazione scientifica}

\begin{esercizio}[\liv{1}]
Scrivere in notazione scientifica: $0{,}000\,32$; \ $45\,000\,000$; \ $0{,}0605$; \ $1{,}0$.
\end{esercizio}
\risp{$3{,}2\cdot10^{-4}$; $4{,}5\cdot10^{7}$; $6{,}05\cdot10^{-2}$; $1{,}0\cdot10^{0}$.}

\begin{esercizio}[\liv{1}]
Calcolare, dando il risultato in notazione scientifica:
\begin{gather*}
(2{,}5\cdot10^{6})(4\cdot10^{-9}),\qquad \frac{6\cdot10^{3}}{1{,}5\cdot10^{-4}},\\
(3\cdot10^{5})^{2},\qquad \sqrt{1{,}6\cdot10^{9}}.
\end{gather*}
\end{esercizio}
\risp{$1\cdot10^{-2}$; \ $4\cdot10^{7}$; \ $9\cdot10^{10}$; \ $\sqrt{16\cdot10^{8}}=4\cdot10^{4}$.}

\begin{esercizio}[\liv{1}]
Indicare il numero di cifre significative di: $0{,}004\,50$; \ $4{,}50\cdot10^{3}$; \ $1{,}200\cdot10^{3}$; \ $0{,}1$. Perch\'e la scrittura ``$1200$'' \`e ambigua?
\end{esercizio}
\risp{$3$; $3$; $4$; $1$.}

\begin{esercizio}[\liv{2}]
Il raggio della Terra \`e circa $6{,}37\cdot10^{6}$ m. Calcolare in notazione scientifica, con tre cifre significative, la superficie $4\pi r^{2}$ e il volume $\frac43\pi r^{3}$ della Terra.
\end{esercizio}
\risp{$5{,}10\cdot10^{14}$ m$^{2}$; \ $1{,}08\cdot10^{21}$ m$^{3}$.}

\begin{esercizio}[\liv{2}]
Un batterio si divide ogni $20$ minuti. Partendo da un solo batterio, e supponendo che tutti si dividano, quanti batteri ci sono dopo $10$ ore? Dare il risultato in notazione scientifica. (Si usi $2^{10}\approx10^{3}$.)
\end{esercizio}
\risp{$2^{30}=(2^{10})^{3}\approx10^{9}$ (il valore esatto \`e $1\,073\,741\,824\approx1{,}07\cdot10^{9}$).}

%-------------------------------------------------------------------
\section{Percentuali}

\begin{esercizio}[\liv{1}]
Calcolare: il $35\,\%$ di $240$; lo $0{,}5\,\%$ di $3000$; la percentuale che $18$ rappresenta rispetto a $150$; il numero di cui $150$ \`e il $60\,\%$.
\end{esercizio}
\risp{$84$; $15$; $12\,\%$; $250$.}

\begin{esercizio}[\liv{1}]
Una popolazione passa da $1250$ a $1100$ individui. Calcolare la variazione assoluta e la variazione relativa.
\end{esercizio}
\risp{$-150$ individui; $-12\,\%$.}

\begin{esercizio}[\liv{2}]
Dopo una diminuzione del $25\,\%$, una popolazione conta $900$ individui. Quanti erano prima della diminuzione?
\end{esercizio}
\risp{$\frac{900}{0{,}75}=1200$ (non $900\cdot1{,}25=1125$).}

\begin{esercizio}[\liv{2}]
La copertura forestale di una provincia passa dal $32\,\%$ al $28\,\%$ del territorio. Esprimere la variazione in punti percentuali e in percentuale.
\end{esercizio}
\risp{$-4$ punti percentuali; $-12{,}5\,\%$.}

%-------------------------------------------------------------------
\section{Problemi di riepilogo}

\begin{esercizio}[\liv{1}]
Un'aula misura $8\times6\times3$ m. Calcolarne il volume in $\mathrm{m^3}$ e in litri, e scrivere il secondo risultato in notazione scientifica.
\end{esercizio}
\risp{$144\ \mathrm{m^3}=144\,000$ L $=1{,}44\cdot10^5$ L.}

\begin{esercizio}[\liv{2}]
Un lago di superficie $3{,}2\ \mathrm{km^2}$ e profondità media $4{,}5$ m viene raccolto, idealmente, in una vasca cubica. Quanto misura il lato? Esprimere il risultato in notazione scientifica con due cifre significative.
\end{esercizio}
\risp{Volume $1{,}44\cdot10^7\ \mathrm{m^3}$; lato $\sqrt[3]{1{,}44\cdot10^7}\ \mathrm m\approx2{,}4\cdot10^2$ m.}

%-------------------------------------------------------------------
\stampasoluzioni
